\scanpage{192}
et les
\[
\boxed{B_i\in\mathcal A_1^0=k\{X_0,X_1\}\qquad(i\ge0)}
\]
par les formules
\[
\begin{cases}
\varphi(X_0^{(i)})=(X_0+f)^{(i)}+B_i\,dX_0\\
\varphi(dX_0)=1+df
\end{cases}
\]
(de sorte qu'on a
\[
f=\varphi_0^0(X_0)-X_0=-(\varphi_0^0(X_1)-X_1)
\]
\[
B_i\,\unsure{dX_0}=\varphi^1(X_0^{(i)})
\]
)

\uline{NB.} Si on échange le rôle de $X_0$ et $X_1$, $f$ change de signe, que deviennent les $B_i$ ? Il faut calculer
\[
\begin{aligned}
\varphi^1(X_1^{(i)})&=\varphi^1((T-X_0)^{(i)})\\
&=\varphi^1\left(\sum_{0\le j\le i}T^{(i-j)}X_0^{(j)}(-1)^{\unsure{j}}\right)
=\left(\sum_{0\le j\le i}T^{(i-j)}(-1)^jB_j\right)dX_0\\
&=\underbrace{\left(\sum_{0\le j\le i}(-1)^{j+1}T^{(i-j)}B_j\right)}_{C_i}dX_1
\end{aligned}
\]
\deleted{\unsure{En termes}}
\deleted{$\varphi^1(T)$}
\[
\underbrace{\varphi^1}_{\unsure{\pi^1\varphi|\operatorname{pr}^1\varphi}}
\left(
\underbrace{(X_0+X_1)^{(i)}}_{
\underbrace{\displaystyle\sum_{j+k=i}X_0^{(j)}X_1^{(k)}}_{
\displaystyle\unsure{\pi^1|\operatorname{pr}^1}\sum_{j+k=i}
\underbrace{\varphi(X_0^{(j)})}_{(X_0+f)^{(j)}+B_jdX_0}
\underbrace{\varphi(X_1^{(k)})}_{(X_1-f)^{(k)}+C_k\underbrace{dX_1}_{-dX_0}}
}}
\deleted{\unsure{+}}
\right)=\varphi^1(T^{\unsure{(i)|i}})=0
\]
donc
\[
\sum_{j+k=i}\left[-(X_0+f)^{(j)}C_k\unsure{+}B_j(X_1-f)^{(k)}\right]=0
\]
\margin{191}
